% -----------------------------------------------------------------------------
% File: sections/01_introduction.tex
% Description: Section 1 - General Introduction
% -----------------------------------------------------------------------------

\section{General Introduction}
\label{sec:intro}

\subsection{Context and Research Motivation}

The rapid convergence of deep learning architectures and large language models (LLMs) has catalyzed a paradigm shift across diverse industrial and scientific domains. Among these, two fields stand out for their critical importance, high standards of accuracy, and unique engineering challenges: \highlight{Computer-Aided Diagnosis (CAD)} in medical imaging and \highlight{Automated Software Engineering (ASE)} in software development.

In medical imaging, particularly in hematological oncology and neurodegenerative disease screening, deep learning has transitioned from an experimental tool to a clinical necessity. The diagnosis of diseases like leukemia and Alzheimer's relies heavily on high-precision analysis of visual data, such as peripheral blood smear images and magnetic resonance imaging (MRI) scans. However, manual examination of these materials is notoriously time-consuming, prone to subjective variance, and constrained by a global shortage of clinical specialists. The main challenges here are visual: distinguishing subtle cellular morphology aberrations and detecting localized structural brain tissue shrinkage (atrophy) from high-dimensional scans.

Concurrently, automated software engineering has entered a new era driven by generative artificial intelligence. Code review, which has historically been a human-centric bottleneck in the software development lifecycle (SDLC), is being automated using LLMs as code review agents. This paradigm promises to dramatically reduce developer overhead, accelerate time-to-market, and enforce consistent code quality standards. Nevertheless, deploying LLM agents in production software environments introduces critical questions regarding scoring consistency, domain alignment, data privacy, and the trade-off between API costs and actual quality improvements.

\subsection{Report Objectives}

This report presents a comprehensive analysis and synthesis of three state-of-the-art research papers presented at the International Conference on Computational Collective Intelligence (ICCCI 2025), addressing the above challenges:
\begin{enumerate}
    \item \highlight{Leukemia Detection Based on YOLOv11 with Global and Local Contexts Interaction} \cite{leukemia_gcli_2025}: proposes integrating a lightweight attention module (GCLI) into the YOLOv11 architecture to improve blast cell detection in peripheral blood smear images.
    \item \highlight{The New Method for Detection of Alzheimer's Disease} \cite{alzheimer_vgg_2025}: explores subject-level MRI classification using progressive transfer learning with VGG16 and a dual-model disjunction fusion strategy.
    \item \highlight{LLMs as Code Review Agents: A Rapid Review and Experimental Evaluation with Human Expert Judges} \cite{llm_review_2025}: conducts a comparative analysis evaluating LLM-generated code quality scores against senior developer consensus across 12 quality metrics.
\end{enumerate}

For each paper, we present the research background, the proposed methodology, the experimental setup and dataset, and a detailed analysis of the results. We conclude each section with remarks on the paper's contributions and limitations.
